\section{System Model and Problem Formulation}

\subsection{Multimodal Observation Model}

We consider \lb{an empirical deadline-aware} multimodal diagnosis system deployed in cloud microservice environments. At each inference step \(t\), the system consumes a sliding-window snapshot of preprocessed multimodal observations:

\begin{equation}
\mathcal{O}_{t-W+1:t} =
\{\mathbf{X}^{m}_{t-W+1:t}, \mathbf{X}^{l}_{t-W+1:t}, \mathbf{X}^{tr}_{t-W+1:t}\},
\end{equation}
where \(\mathbf{X}^{m}\in\mathbb{R}^{W\times N\times d_m}\) denotes service-level metric features, \(\mathbf{X}^{l}\in\mathbb{R}^{W\times N\times d_l}\) denotes log-derived statistical features, and \(\mathbf{X}^{tr}\in\mathbb{R}^{W\times N\times N\times d_{tr}}\) denotes trace edge features over service pairs. Together, they capture complementary facets of system behavior: metrics reflect resource usage patterns such as CPU, memory, and network traffic, logs summarize the distribution of runtime events such as error and warning levels, and traces encode request-flow relations together with edge statistics such as call counts and durations. The joint window \(\mathcal{O}_{t-W+1:t}\) thus serves as the input to the online diagnosis process.

\subsection{Diagnosis Objective}

Given the observation window \(\mathcal{O}_{t-W+1:t}\), the diagnosis system aims to produce an online anomaly decision together with service-level ranking and attribution signals. Specifically, the system determines whether the current window is abnormal, assigns anomaly scores to services, and exposes diagnostic evidence that can support subsequent root-cause analysis by operators or downstream tools.

These tasks must be performed in an online manner under a predefined latency budget $D$, which is determined by the response requirements of the target cloud application. \lb{We treat this as an empirical soft real-time requirement evaluated through online replay: deadline misses and observed tail latency are measured directly, while formal WCET and schedulability bounds are outside the scope of the current prototype.}
\subsection{Inference and Latency Model}

We model the diagnosis process as a parametric inference function \(f\) that maps the observation window \(\mathcal{O}_{t-W+1:t}\) to a predicted diagnosis result \(\hat{y}(t)\):
\[
\hat{y}(t) = f(\mathcal{O}_{t-W+1:t}; k(t)),
\]
where \lb{\(k(t)\) denotes the runtime sparse expert budget, i.e., the number of experts activated by the router at each adapted layer under a bounded maximum budget selected by a validation-fixed confidence-threshold policy}. In practice, latency is affected by fixed preprocessing, device transfer, kernel scheduling, and sparse expert execution. We therefore use the following abstract decomposition:
\[
T(t) = T_{\mathrm{base}}(t) + T_{\mathrm{sparse}}(t;k(t)),
\]
where \(T_{\mathrm{base}}(t)\) includes data preparation and backbone execution, and \lb{\(T_{\mathrm{sparse}}(t;k(t))\) is the additional cost of sparse expert adaptation under the selected runtime budget}. The system aims to keep the observed response time within the deadline budget
\[
T(t) \leq D,
\]
for every inference step \(t\).


\subsection{Adaptive Inference Control}

\lb{The sparse budget \(k(t)\) determines how many experts are evaluated for the current routed representation, subject to a maximum budget \(k_{\max}\). ASID learns router scores and service-conditioned prior weights during training, while the hard budget threshold is selected on the validation split and then fixed before test replay. At inference time, the confidence-threshold policy maps the router top-1 confidence to \(k(t)\), and the router selects the expert indices under that budget. Larger budgets may increase representational capacity but also raise computational cost; smaller budgets reduce cost but may weaken expert specialization.}

\subsection{Problem Formulation}

Our objective is to design a sparse inference policy that maximizes diagnostic performance under latency constraints.

\begin{align}
\lb{\max_{\theta,\,\rho} \quad} & \lb{\mathbb{E}[A(t)]} \\
\text{s.t.} \quad
& \lb{T(t) \leq D,} \\
& \lb{|\rho(\mathcal{O}_{t-W+1:t})| = k(t) \leq k_{\max}.}
\end{align}

\lb{Here \(A(t)\) denotes the diagnosis utility, \(\theta\) denotes trainable model parameters, and \(\rho(\cdot)\) denotes the sparse routing policy. The runtime budget \(k(t)\) is produced by the confidence-adaptive budget rule, and the router then selects the active expert set under that budget.}
